% ======================================================================
%   文件名: main.tex
%   功能: computer-notes.cls 完整示例与测试文档
%   编译方式: xelatex main.tex (或 lualatex main.tex)
%   字符编码: UTF-8
% ======================================================================
\documentclass{computer-notes}

% --- 文档元信息设置 ---
\settitle{计算机科学笔记}
\setauthor{By Your Name}
\setdate{\today}
% ======================================================================
%   文档正文开始
% ======================================================================


\begin{document}

% 生成封面
\makecover

% 生成目录
\tableofcontents
\newpage

\part{Part 测试}


% ======================================================================
%   第一章: 标题与正文演示
% ======================================================================

\chapter{测试}

\section{标题演示与正文格式}

这是 \textbf{section} 级别的标题。\emphasis{本节演示文档的基本排版格式设置。}


计算机科学是一门理论与实践相结合的学科。学习编程需要掌握算法、数据结构、设计模式等核心概念。同时，实践能力的培养同样重要——通过大量的项目实践，我们可以深化对理论知识的理解。

\ssection{三级标题示例}

这是 \textbf{ssection} 级别的标题。在实际编程中，代码的可读性至关重要。良好的命名规范、适当的注释和清晰的代码结构可以大大提升代码的可维护性。


\sssction{四级标题示例}

这是 \textbf{sssction} 级别的标题。四级标题通常用于更细粒度的内容分组。

\paragraph{五级标题示例}

这是 \textbf{paragraph} 级别的标题。适合非常细的小节说明。

\subparagraph{六级标题示例}

这是 \textbf{subparagraph} 级别的标题。适合补充型说明内容。

% ======================================================================
%   列表演示
% ======================================================================


\section{列表与排版}

\ssection{无序列表}


编程的基本要素包括：

\begin{itemize}
  \item 变量与数据类型
  \item 控制流（if-else, loops）
  \item 函数与模块化
  \item 面向对象编程
  \item 设计模式
  \item 算法与复杂度分析
\end{itemize}

\ssection{有序列表}

学习编程的推荐步骤：

\begin{enumerate}
  \item 学习基本语法和数据类型
  \item 理解控制流和函数
  \item 深入学习数据结构
  \item 掌握常见算法
  \item 学习面向对象编程
  \item 应用设计模式
  \item 进行大规模项目实践
\end{enumerate}

\ssection{嵌套列表}

数据结构分类：

\begin{itemize}
  \item \textbf{线性结构}
    \begin{itemize}
      \item 数组
      \item 链表
      \item 栈
      \item 队列
    \end{itemize}
  \item \textbf{树形结构}
    \begin{itemize}
      \item 二叉树
      \item 平衡树
      \item 堆
    \end{itemize}
  \item \textbf{图形结构}
    \begin{itemize}
      \item 有向图
      \item 无向图
      \item 加权图
    \end{itemize}
\end{itemize}

% ======================================================================
%   数学公式演示
% ======================================================================

\section{数学公式}

\ssection{行内公式}

复杂度分析中，我们常用大 O 记号。例如，线性搜索的时间复杂度为 $O(n)$，二分搜索为 $O(\log n)$。

\ssection{方程式环境}

二次方程的求根公式为：

\begin{equation}
  x = \frac{-b \pm \sqrt{b^2 - 4ac}}{2a}
\end{equation}

\ssection{多行公式}

斐波那契数列的递推关系：

\begin{align}
  F(0) &= 0 \\
  F(1) &= 1 \\
  F(n) &= F(n-1) + F(n-2), \quad n \geq 2
\end{align}

\ssection{非编号公式}

快速排序的分治策略为：

\begin{equation*}
  T(n) = \begin{cases}
    O(1) & \text{if } n \leq 1 \\
    T(q) + T(n-q-1) + O(n) & \text{if } n > 1
  \end{cases}
\end{equation*}


% ======================================================================
%   Callout 提示框演示 (7 种类型)
% ======================================================================
\chapter{测试2 - English Text} 
\section{提示框系统 (Callout)}

\ssection{笔记框}

\note{%
这是一个笔记框示例。使用 \texttt{\textbackslash{}note\{内容\}} 命令可以快速插入笔记。\\
笔记框通常用于补充说明或重要备注。
}

\ssection{技巧框}

\tip{%
这是一个技巧框。使用 \texttt{\textbackslash{}tip\{内容\}} 命令。\\
技巧框用于分享编程技巧或最佳实践。例如：使用 \texttt{const} 声明不可变变量，有利于编写更安全的代码。
}

\ssection{信息框}

\info{%
这是一个信息框。使用 \texttt{\textbackslash{}info\{内容\}} 命令。\\
信息框用于传达一般性信息或背景知识。\\
例如：JavaScript 最初设计用于浏览器端的脚本编程，现在已广泛用于服务器端和桌面应用。
}

\ssection{警告框}

\warning{%
这是一个警告框。使用 \texttt{\textbackslash{}warning\{内容\}} 命令。\\
警告框用于提醒注意某些可能的问题。例如：\textit{在使用指针时要小心内存泄漏问题。}
}

\ssection{注意框}

\caution{%
这是一个注意框。使用 \texttt{\textbackslash{}caution\{内容\}} 命令。\\
注意框用于强调需要特别关注的事项。例如：修改全局变量可能导致程序状态混乱。
}

\ssection{危险框}

\danger{%
这是一个危险框。使用 \texttt{\textbackslash{}danger\{内容\}} 命令。\\
危险框用于警告严重错误或不应采取的行动。例如：\textbf{不要在生产环境中禁用安全检查！}
}

\ssection{待办框}

\todo{%
这是一个待办框。使用 \texttt{\textbackslash{}todo\{内容\}} 命令。\\
待办框用于记录需要完成的任务。例如：实现异常处理机制。
}

% ======================================================================
%   代码块演示
% ======================================================================

\section{代码块系统}

\ssection{Python 代码块}

下面是一个 Python 函数示例，展示递归计算斐波那契数列：


\begin{codeblock}[Python]{fibonacci.py}
def fibonacci(n):
    """计算第 n 个斐波那契数"""
    if n <= 1:
        return n
    return fibonacci(n - 1) + fibonacci(n - 2)

# 调用示例
result = fibonacci(10)
print(f"第 10 个斐波那契数: {result}")
\end{codeblock}

\tip{%
上面的递归实现虽然直观，但效率较低（时间复杂度 $O(2^n)$）。实际应用中应使用动态规划或迭代方式优化。
}

\ssection{Bash 脚本示例}

\begin{codeblock}[Bash]{backup.sh}
#!/bin/bash
# 简单的备份脚本

SOURCE_DIR="/home/user/documents"
BACKUP_DIR="/mnt/backup"
DATE=$(date +%Y%m%d_%H%M%S)

cp -r "$SOURCE_DIR" "$BACKUP_DIR/backup_$DATE"
echo "备份完成: $BACKUP_DIR/backup_$DATE"
\end{codeblock}

\ssection{行内代码示例}

当我们需要在正文中插入代码片段时，可以使用 \inline{inline code} 命令。例如，定义变量时使用 \inline{const x = 10;} 比 \inline{var x = 10;} 更安全。

% ======================================================================
%   其他组件演示
% ======================================================================

\section{其他实用组件}

\ssection{行内代码与强调}

编程中常见的特殊值包括:
\begin{itemize}
  \item \inline{null} 和 \inline{undefined}
  \item \inline{true} 和 \inline{false}
  \item \inline{NaN} (Not a Number)
\end{itemize}

强调内容: \emphasis{始终使用类型检查来避免类型转换的陷阱}。

\ssection{定义框}

\definition{算法}{%
算法是一步步的过程描述，用于解决某个具体问题。一个好的算法应该具备：正确性、有限性、输入/输出和有效性。
}

\definition{数据结构}{%
数据结构是计算机中组织与存储数据的方式。选择合适的数据结构对提高程序性能至关重要。常见的数据结构包括数组、链表、树、图等。
}

\ssection{算法框}

\algorithm{%
\textbf{排序算法 - 冒泡排序}

\begin{enumerate}
  \item 遍历数组，比较相邻的两个元素
  \item 若左元素 > 右元素，交换位置
  \item 重复该过程直到数组有序
  \item 时间复杂度：最坏 $O(n^2)$，最好 $O(n)$
\end{enumerate}
}


\ssection{快捷键样式}

保存文件: \shortcut{Ctrl + S}

打开终端: \shortcut{Ctrl + `}

全选: \shortcut{Ctrl + A}

复制: \shortcut{Ctrl + C}

粘贴: \shortcut{Ctrl + V}

\ssection{终端命令}

查看文件内容:

\terminal{cat file.txt}

列出目录内容:

\terminal{ls -la /home/user}

创建新文件:

\terminal{touch myfile.txt}

\ssection{文件路径示例}

配置文件通常位于 \file{~/.config/app/config.yml}。

系统日志在 \file{/var/log/system.log}。

用户目录是 \file{/home/username}。

\ssection{文本高亮}

这是 \highlight{高亮显示的内容}。适用于需要特别强调的词汇或短语。

% ======================================================================
%   综合应用示例
% ======================================================================

\section{综合应用: Web 开发基础}

\ssection{概述}

Web 开发涉及前端、后端和数据库等多个层面。学习 Web 开发需要掌握 HTML、CSS、JavaScript 等基础技术。

\ssection{前端技术栈}

\note{%
现代前端开发已不再是简单的 HTML + CSS + JavaScript 的组合。React、Vue、Angular 等框架的出现让前端开发变得更加结构化和高效。
}

前端常用技术：

\begin{itemize}
  \item \textbf{HTML5}: 结构与语义化标记
  \item \textbf{CSS3}: 样式与布局
  \item \textbf{JavaScript}: 交互逻辑
  \item \textbf{前端框架}: React/Vue/Angular
  \item \textbf{构建工具}: Webpack/Vite
\end{itemize}

\ssection{后端技术栈}

\info{%
后端的核心职责是处理业务逻辑、数据存储与验证。选择合适的语言和框架应根据项目需求而定。
}

后端常用技术：

\begin{enumerate}
  \item Node.js / Express (JavaScript)
  \item Python / Django / Flask
  \item Java / Spring Boot
  \item Go / Gin
  \item 数据库: MySQL / PostgreSQL / MongoDB
\end{enumerate}

\ssection{REST API 设计}

RESTful API 设计的基本原则：

\begin{itemize}
  \item GET: 获取资源
  \item POST: 创建资源
  \item PUT: 更新资源
  \item DELETE: 删除资源
  \item PATCH: 部分更新
\end{itemize}

API 响应示例:

\begin{codeblock}[Python]{api\_response.py}
# 成功响应
response = {
    "status": 200,
    "code": "SUCCESS",
    "data": {
        "id": 1,
        "name": "用户名",
        "email": "user@example.com"
    },
    "message": "操作成功"
}

# 错误响应
error_response = {
    "status": 400,
    "code": "INVALID_INPUT",
    "message": "参数验证失败",
    "errors": [
        {"field": "email", "message": "邮箱格式不正确"}
    ]
}
  \end{codeblock}

\warning{%
API 设计中要特别注意安全问题：
\begin{itemize}
  \item 实施身份认证 (JWT / OAuth2)
  \item 进行速率限制 (Rate Limiting)
  \item 验证用户输入 (Input Validation)
  \item 使用 HTTPS 加密传输
\end{itemize}
}

\ssection{数据库基础}

\definition{数据库}{%
数据库是结构化数据的集合，通过数据库管理系统 (DBMS) 进行管理。关系数据库使用表格结构存储数据，并通过 SQL 语言进行查询。
}

关键概念：

\begin{itemize}
  \item \textbf{表}: 行和列的集合
  \item \textbf{主键}: 唯一标识记录的字段
  \item \textbf{外键}: 建立两个表之间的关系
  \item \textbf{索引}: 加快查询速度的数据结构
  \item \textbf{约束}: 确保数据完整性的规则
\end{itemize}

基础 SQL 操作：

\begin{codeblock}[SQL]{sql\_examples.sql}
-- 创建表
CREATE TABLE users (
    id INT PRIMARY KEY AUTO_INCREMENT,
    name VARCHAR(100) NOT NULL,
    email VARCHAR(100) UNIQUE,
    created_at TIMESTAMP DEFAULT CURRENT_TIMESTAMP
);

-- 插入数据
INSERT INTO users (name, email) VALUES ('Alice', 'alice@example.com');

-- 查询数据
SELECT * FROM users WHERE id = 1;

-- 更新数据
UPDATE users SET email = 'newemail@example.com' WHERE id = 1;

-- 删除数据
DELETE FROM users WHERE id = 1;
\end{codeblock}

\tip{%
使用 ORM (Object-Relational Mapping) 库如 SQLAlchemy (Python)、Sequelize (Node.js) 可以简化数据库操作，提高代码可读性。
}

% ======================================================================
%   性能优化章节
% ======================================================================

\section{性能优化与最佳实践}

\ssection{算法优化}

\definition{时间复杂度}{%
时间复杂度描述算法运行时间与输入规模的关系。通常使用 Big O 记号表示。
}

常见的时间复杂度从低到高排列：

\begin{equation*}
O(1) < O(\log n) < O(n) < O(n \log n) < O(n^2) < O(n^3) < O(2^n) < O(n!)
\end{equation*}

\caution{%
避免使用低效的算法。例如，对大数据集的嵌套循环可能导致程序性能严重下降。应该优先选择更优的算法。
}

\ssection{代码优化技巧}

\tip{%
常见的代码优化技巧：
\begin{enumerate}
  \item 避免不必要的对象创建
  \item 使用缓存减少重复计算
  \item 选择合适的数据结构
  \item 避免深度递归（可能导致栈溢出）
  \item 使用异步编程处理 I/O 操作
\end{enumerate}
}

记忆化 (Memoization) 优化斐波那契数列：

\begin{codeblock}[Python]{fibonacci\_memo.py}
# 使用字典缓存已计算的结果
def fibonacci_memo(n, memo={}):
    if n in memo:
        return memo[n]
    if n <= 1:
        return n
    result = fibonacci_memo(n - 1, memo) + fibonacci_memo(n - 2, memo)
    memo[n] = result
    return result

# 时间复杂度降低到 O(n)
print(fibonacci_memo(100))
\end{codeblock}

\danger{%
不要过度优化。优化应该基于实际的性能瓶颈分析，而不是盲目猜测。使用分析工具 (Profiler) 来找出真正的性能问题。
}

\ssection{内存管理}

\info{%
在 Python 等自动内存管理语言中，垃圾回收会自动释放不再使用的对象。但在 C/C++ 等手动管理的语言中，必须小心处理内存申请与释放。
}

内存泄漏的常见原因：

\begin{itemize}
  \item 忘记释放动态分配的内存
  \item 循环引用导致无法被垃圾回收
  \item 监听器未注销
  \item 缓存不适当的增长
\end{itemize}

% ======================================================================
%   文档结尾
% ======================================================================

\section{总结与展望}

本笔记涵盖了计算机科学的基础知识，包括编程基础、数据结构、算法、Web 开发等内容。

\definition{学习路径}{%
一个系统的学习路径应该包括：
\begin{enumerate}
  \item 掌握编程基础与常用算法
  \item 学习数据结构与复杂度分析
  \item 深入学习特定领域 (Web、移动、AI 等)
  \item 进行实战项目并不断积累经验
  \item 阅读他人优秀代码，学习最佳实践
\end{enumerate}
}

\todo{%
持续学习的计划：
\begin{itemize}
  \item 每周阅读一篇技术文章
  \item 每月完成一个小项目
  \item 定期回顾基础知识
  \item 参与开源项目贡献
\end{itemize}
}

感谢阅读本笔记！祝您学习愉快！

\end{document}

% ======================================================================
% main.tex 结束
% ======================================================================
